\section{Related Work}

\subsection{Model-Heterogeneous Federated Learning}

Model heterogeneity breaks the direct parameter-averaging assumption used by standard homogeneous FL. FedDF addresses this mismatch through ensemble distillation over unlabeled data \citep{lin2020feddf}. RHFL considers model-heterogeneous clients with noisy labels and confidence-weighted collaboration \citep{fang2022rhfl}. AugHFL explicitly studies common corruption under model heterogeneity, while RAHFL extends this line with diversity-enhanced local representation learning and asymmetric heterogeneous collaboration \citep{fang2023aughfl,fang2025rahfl}. These methods are the closest technical background for our HFL training pipeline. Our question is different: whether client-specific class--corruption associations make corruption operators directional class cues, and where that shortcut forms.

\subsection{Corruption Robustness and Consistency Learning}

CIFAR-C and ImageNet-C established standard common-corruption evaluations across noise, blur, weather, and digital distortions \citep{hendrycks2019corruptions}. AugMix combines diverse stochastic augmentations with Jensen--Shannon consistency to improve corruption robustness and uncertainty \citep{cubuk2020augmix}. CLE-HFL differs in statistical structure: corruption is deliberately associated with class during client training. Accordingly, our diagnosis compares \emph{different operators for the same semantic source} rather than only checking consistency around one corrupted view.

\subsection{Spurious Correlation, Group Robustness, and Environment Inference}

Shortcut learning and group distribution shift have a broad literature \citep{geirhos2020shortcut}. GroupDRO optimizes worst-group risk with known groups \citep{sagawa2020groupdro}; GEORGE recovers hidden subclasses before group-robust learning \citep{sohoni2020george}; EIIL infers environments without a predefined taxonomy \citep{creager2021eiil}; and JTT upweights examples misclassified by an initial ERM model \citep{liu2021jtt}. CVaR and related DRO objectives emphasize upper-tail loss rather than an explicit class--environment support model \citep{rockafellar2000cvar,duchi2021uniform}. These works constrain our novelty claim: pseudo-environment discovery, difficult-example reweighting, and group robustness are not new. PEW is instead a taxonomy-assisted corruption-family witness, while BER changes the total effective risk mass assigned to class--pseudo-environment supports. Our CVaR-DRO implementation is a protocol-matched empirical tail-risk control, not a line-by-line reproduction of either cited method.

\subsection{Federated Spurious Learning and Counterfactual Diagnosis}

Federated spurious features have already been studied. Personalized federated learning can exploit client-specific spurious features \citep{wang2024pflspurious}; FedPIN targets shortcut-averse personalized invariant learning \citep{tang2024fedpin}; and FedCD studies spurious correlation in federated domain generalization \citep{ma2024fedcd}. In the current citation audit, FedCD is treated only as an arXiv preprint. Counterfactual shortcut evaluation also has precedents, including medical-imaging analyses that test whether acquisition or demographic attributes are actually used by a model \citep{vigneshwaran2026counterfactual,jabbour2020chestxray}. Our contribution is not counterfactual analysis in the abstract. We specialize it to source-paired corruption-operator interventions, a pre-registered client-specific binding direction, probability-mass alignment, and a matched local-versus-federated formation analysis.
